% ============================================================
% Chapter 2: State of the Art
%   Section: Constructive and Solutions Journalism
% Included from chapters/02-state-of-the-art/state-of-the-art.tex
% ============================================================

\section{Constructive and Solutions Journalism}
\label{sec:constructive-journalism}

Constructive journalism is defined by McIntyre and Gyldensted as
journalism that applies techniques from positive psychology to news
work in order to create more productive, engaging stories while
holding to journalism's core functions \cite{mcintyre2017constructive}.
It does not mean only good news: a constructive story reports a problem
together with the responses to it, the evidence that they work, and
what readers can take from it. Solutions journalism, a close relative,
centres the story on a response to a problem and on the evidence of its
results. Both react to a news supply that leans to the negative, which
audiences attend to more strongly \cite{soroka2015negativity} and which
can leave readers informed but discouraged.

For this work the publication's \emph{mission} is that kind of news on
five groups of topics: society, science, environment, culture and
health. In practice the mission has to be stated operationally, because
the same outlets publish both kinds of story. The system's discovery
stage keeps what a constructive publication would offer and leaves out
what it would not --- party politics, crime, war, disasters, celebrity,
sport, markets, service pages --- and the line between hard news and
constructive news on the same topic is left to a judgement, by a person
or a language model (Section~\ref{sec:design-selection}).

\paragraph{Why its claims are harder to check.} Fact-checking practice
and its datasets were built around political statements: a quote, a
speaker, a disputed figure. The claims of constructive news are
different in ways that make them harder, not easier, to verify
automatically.

\begin{itemize}
    \item \textbf{They are new.} A study published this week or a
    project launched this month has little independent coverage yet,
    and the open web may hold nothing but the announcement itself.
    \item \textbf{They are often sourced to their protagonist.} The
    figures of a positive story frequently come from the press release
    of the organisation the story is about. That confirms what was
    said, not that it is true, and the annotation guide of this work
    treats such a claim as not enough evidence
    (Section~\ref{subsec:custom-rules}).
    \item \textbf{They mix facts and interpretation.} ``The first
    hospital in the region to\ldots'' is checkable; ``a hopeful step
    for the sector'' is not, and both appear in the same paragraph.
    \item \textbf{They are rarely false.} The custom validation set of
    this work found no refuted claim among 109 natural ones
    (Section~\ref{subsec:custom-size}). Errors, when they occur, are in
    a detail --- a figure, a date, an attribution --- which is the
    \texttt{PARTIALLY\_TRUE} case rather than the \texttt{FALSE} one.
\end{itemize}
